\documentclass[10pt,a4paper]{amsart}
\usepackage{amsmath,amssymb,amsthm,graphicx}
\usepackage{verbatim}
\newenvironment{reference}{\comment}{\endcomment}
%\newenvironment{reference}{}{}
\newenvironment{slogan}{\comment}{\endcomment}
\newenvironment{history}{\comment}{\endcomment}

% For commutative diagrams we use Xy-pic
\usepackage[all]{xy}

% We use 2cell for 2-commutative diagrams.
\xyoption{2cell}
% Theorem environments.
%
\theoremstyle{plain}
\newtheorem{theorem}[subsection]{Theorem}
\newtheorem{proposition}[subsection]{Proposition}
\newtheorem{lemma}[subsection]{Lemma}

\theoremstyle{definition}
\newtheorem{definition}[subsection]{Definition}
\newtheorem{example}[subsection]{Example}
\newtheorem{exercise}[subsection]{Exercise}
\newtheorem{situation}[subsection]{Situation}

\theoremstyle{remark}
\newtheorem{remark}[subsection]{Remark}
\newtheorem{remarks}[subsection]{Remarks}

\numberwithin{equation}{subsection}

% Macros
%
\def\lim{\mathop{\mathrm{lim}}\nolimits}
\def\colim{\mathop{\mathrm{colim}}\nolimits}
\def\Spec{\mathop{\mathrm{Spec}}}
\def\Hom{\mathop{\mathrm{Hom}}\nolimits}
\def\Ext{\mathop{\mathrm{Ext}}\nolimits}
\def\SheafHom{\mathop{\mathcal{H}\!\mathit{om}}\nolimits}
\def\SheafExt{\mathop{\mathcal{E}\!\mathit{xt}}\nolimits}
\def\Sch{\mathit{Sch}}
\def\Mor{\mathop{\mathrm{Mor}}\nolimits}
\def\Ob{\mathop{\mathrm{Ob}}\nolimits}
\def\Sh{\mathop{\mathit{Sh}}\nolimits}
\def\NL{\mathop{N\!L}\nolimits}
\def\CH{\mathop{\mathrm{CH}}\nolimits}
\def\proetale{{pro\text{-}\acute{e}tale}}
\def\etale{{\acute{e}tale}}
\def\QCoh{\mathit{QCoh}}
\def\Ker{\mathop{\mathrm{Ker}}}
\def\Im{\mathop{\mathrm{Im}}}
\def\Coker{\mathop{\mathrm{Coker}}}
\def\Coim{\mathop{\mathrm{Coim}}}

% Boxtimes
%
\DeclareMathSymbol{\boxtimes}{\mathbin}{AMSa}{"02}

%
% Macros for moduli stacks/spaces
%
\def\QCohstack{\mathcal{QC}\!\mathit{oh}}
\def\Cohstack{\mathcal{C}\!\mathit{oh}}
\def\Spacesstack{\mathcal{S}\!\mathit{paces}}
\def\Quotfunctor{\mathrm{Quot}}
\def\Hilbfunctor{\mathrm{Hilb}}
\def\Curvesstack{\mathcal{C}\!\mathit{urves}}
\def\Polarizedstack{\mathcal{P}\!\mathit{olarized}}
\def\Complexesstack{\mathcal{C}\!\mathit{omplexes}}
% \Pic is the operator that assigns to X its picard group, usage \Pic(X)
% \Picardstack_{X/B} denotes the Picard stack of X over B
% \Picardfunctor_{X/B} denotes the Picard functor of X over B
\def\Pic{\mathop{\mathrm{Pic}}\nolimits}
\def\Picardstack{\mathcal{P}\!\mathit{ic}}
\def\Picardfunctor{\mathrm{Pic}}
\def\Deformationcategory{\mathcal{D}\!\mathit{ef}}
\usepackage{cleveref}
\crefname{originallemma}{original lemma}{original lemmas}
\Crefname{originallemma}{Original lemma}{Original lemmas}
\makeatletter
\newlabel{more-morphisms-lemma-Noetherian-pseudo-coherent}{{0684}{source}{}{sourcepointer-1}{}}
\newlabel{more-morphisms-lemma-Noetherian-pseudo-coherent@cref}{{[originallemma][0][]0684}{[1][0][]source}{}{}{}}
\makeatother
\crefname{originallemma}{original lemma}{original lemmas}
\Crefname{originallemma}{Original lemma}{Original lemmas}
\makeatletter
\newlabel{morphisms-lemma-i-star-equivalence}{{01QY}{source}{}{sourcepointer-2}{}}
\newlabel{morphisms-lemma-i-star-equivalence@cref}{{[originallemma][0][]01QY}{[1][0][]source}{}{}{}}
\makeatother
\crefname{originallemma}{original lemma}{original lemmas}
\Crefname{originallemma}{Original lemma}{Original lemmas}
\makeatletter
\newlabel{lemma-sheaf-with-exact-support-ext}{{0A77}{source}{}{sourcepointer-3}{}}
\newlabel{lemma-sheaf-with-exact-support-ext@cref}{{[originallemma][0][]0A77}{[1][0][]source}{}{}{}}
\makeatother
\crefname{originallemma}{original lemma}{original lemmas}
\Crefname{originallemma}{Original lemma}{Original lemmas}
\makeatletter
\newlabel{perfect-lemma-quasi-coherence-internal-hom}{{0A6H}{source}{}{sourcepointer-4}{}}
\newlabel{perfect-lemma-quasi-coherence-internal-hom@cref}{{[originallemma][0][]0A6H}{[1][0][]source}{}{}{}}
\makeatother
\crefname{originallemma}{original lemma}{original lemmas}
\Crefname{originallemma}{Original lemma}{Original lemmas}
\makeatletter
\newlabel{lemma-compute-sheaf-with-exact-support}{{0A75}{source}{}{sourcepointer-5}{}}
\newlabel{lemma-compute-sheaf-with-exact-support@cref}{{[originallemma][0][]0A75}{[1][0][]source}{}{}{}}
\makeatother
\title{The sheaf with exact support preserves quasi-coherent complexes and agrees with the algebraic Hom on affine charts}
\author{The Stacks Project Authors}
\date{stacks-project:duality.tex:local-snapshot-20261012; GNU FDL 1.2 or later}
\begin{document}
\maketitle
\noindent\textit{Editorial note.} All delivered mathematical text is the author's own, copied from the pinned Stacks Project chapter file duality.tex; no statement or proof step was rewritten, reordered or removed.\par\medskip
\noindent\textit{Editorial note.} Preamble provenance: the pinned chapter file opens with an input of the shared preamble and carries no preamble text of its own, so the project's single Stacks-wide preamble staged at sources/stacks-duality/preamble.tex is excerpted verbatim under the recipe fixed by LANE\_BRIEF 22.1, as three spans, lines 8-13 (verbatim package and the reference/slogan/history environments), lines 14-18 (Xy-pic setup) and lines 154-217 (the native theorem declarations and mathematical macros). Nothing was added, removed or redefined; the wrapper class and the amsmath/amssymb/amsthm/graphicx base packages are supplied by the assembler.\par\medskip
\noindent\textit{Editorial note.} The counted statement and proof contain no cross reference, no citation and no included file, so the delivered body is self contained.\par\medskip
\noindent\textit{Editorial note.} Canonical location of the counted proposition: Stacks tag 0A78 (sources/stacks-duality/pinned-original-source.tex line 2373).\par\medskip
\noindent\textit{Allowed standard prerequisites (editorial).}\begin{itemize}
\item The sheaf with exact support and its computation. Same chapter, Lemma labelled lemma-sheaf-with-exact-support-ext, cited at proof lines 2400-2404 and 2438-2442, and Lemma labelled lemma-compute-sheaf-with-exact-support, cited at proof lines 2424-2428. Identifies RSheafHom(O\_Z, -) with the object whose quasi-coherence is checked.
\item Reflecting quasi-coherence along the pushforward of a closed immersion. Morphisms, Lemma labelled morphisms-lemma-i-star-equivalence, cited at proof lines 2398-2400. Reduces (1) to a statement about i\_* of the cohomology sheaves.
\item Internal Hom of a pseudo-coherent module stays quasi-coherent. Derived Categories of Schemes, Lemma labelled perfect-lemma-quasi-coherence-internal-hom, cited at proof lines 2408-2412 and 2442-2446. Gives (1) and identifies the affine comparison map.
\item Pseudo-coherence of a closed immersion into a locally Noetherian scheme. More on Morphisms, Lemma labelled more-morphisms-lemma-Noetherian-pseudo-coherent, cited at proof lines 2394-2396. Justifies the parenthetical remark in the statement.
\end{itemize}\medskip
\ifdefined\hypertarget\hypertarget{sourcepointer-1}{}\fi
\noindent\textit{Source reference.} Original lemma 0684: sources/stacks-chow/tags.txt, source lines 7993--7993. Cross-chapter standard prerequisite located in the project tag list staged in the local snapshot; invoked at proof line 2395 of the counted proof. It is retained only as a navigation pointer; every load-bearing step is inside the delivered span.\par\medskip
\ifdefined\hypertarget\hypertarget{sourcepointer-2}{}\fi
\noindent\textit{Source reference.} Original lemma 01QY: sources/stacks-chow/tags.txt, source lines 2283--2283. Cross-chapter standard prerequisite located in the project tag list staged in the local snapshot; invoked at proof line 2399 of the counted proof. It is retained only as a navigation pointer; every load-bearing step is inside the delivered span.\par\medskip
\ifdefined\hypertarget\hypertarget{sourcepointer-3}{}\fi
\noindent\textit{Source reference.} Original lemma 0A77: sources/stacks-duality/pinned-original-source.tex, source lines 2339--2339. Same-chapter pointer into the author environment lemma at source line 2339; invoked at proof line 2403 of the counted proof. It is retained only as a navigation pointer; every load-bearing step is inside the delivered span.\par\medskip
\ifdefined\hypertarget\hypertarget{sourcepointer-4}{}\fi
\noindent\textit{Source reference.} Original lemma 0A6H: sources/stacks-chow/tags.txt, source lines 12845--12845. Cross-chapter standard prerequisite located in the project tag list staged in the local snapshot; invoked at proof line 2410 of the counted proof. It is retained only as a navigation pointer; every load-bearing step is inside the delivered span.\par\medskip
\ifdefined\hypertarget\hypertarget{sourcepointer-5}{}\fi
\noindent\textit{Source reference.} Original lemma 0A75: sources/stacks-duality/pinned-original-source.tex, source lines 2283--2283. Same-chapter pointer into the author environment lemma at source line 2283; invoked at proof line 2424 of the counted proof. It is retained only as a navigation pointer; every load-bearing step is inside the delivered span.\par\medskip
\begin{lemma}
\label{lemma-sheaf-with-exact-support-quasi-coherent}
Let $i : Z \to X$ be a pseudo-coherent closed immersion of schemes
(any closed immersion if $X$ is locally Noetherian).
Then
\begin{enumerate}
\item $R\SheafHom(\mathcal{O}_Z, -)$ maps $D^+_\QCoh(\mathcal{O}_X)$
into $D^+_\QCoh(\mathcal{O}_Z)$, and
\item if $X = \Spec(A)$ and $Z = \Spec(B)$, then the diagram
$$
\xymatrix{
D^+(B) \ar[r] & D_\QCoh^+(\mathcal{O}_Z) \\
D^+(A) \ar[r] \ar[u]^{R\Hom(B, -)} &
D_\QCoh^+(\mathcal{O}_X) \ar[u]_{R\SheafHom(\mathcal{O}_Z, -)}
}
$$
is commutative.
\end{enumerate}
\end{lemma}

\begin{proof}
To explain the parenthetical remark, if $X$ is locally Noetherian, then
$i$ is pseudo-coherent by
More on Morphisms, Lemma \ref{more-morphisms-lemma-Noetherian-pseudo-coherent}.

\medskip\noindent
Let $K$ be an object of $D^+_\QCoh(\mathcal{O}_X)$. To prove (1), by
Morphisms, Lemma \ref{morphisms-lemma-i-star-equivalence}
it suffices to show that $i_*$ applied to
$H^n(R\SheafHom(\mathcal{O}_Z, K))$ produces a
quasi-coherent module on $X$. By
Lemma \ref{lemma-sheaf-with-exact-support-ext}
this means we have to show that
$R\SheafHom_{\mathcal{O}_X}(i_*\mathcal{O}_Z, K)$
is in $D_\QCoh(\mathcal{O}_X)$. Since $i$ is pseudo-coherent
the sheaf $\mathcal{O}_Z$ is a pseudo-coherent $\mathcal{O}_X$-module.
Hence the result follows from
Derived Categories of Schemes, Lemma
\ref{perfect-lemma-quasi-coherence-internal-hom}.

\medskip\noindent
Assume $X = \Spec(A)$ and $Z = \Spec(B)$ as in (2).
Let $I^\bullet$ be a bounded below complex of injective $A$-modules
representing an object $K$ of $D^+(A)$.
Then we know that $R\Hom(B, K) = \Hom_A(B, I^\bullet)$ viewed
as a complex of $B$-modules. Choose a quasi-isomorphism
$$
\widetilde{I^\bullet} \longrightarrow \mathcal{I}^\bullet
$$
where $\mathcal{I}^\bullet$ is a bounded below complex of injective
$\mathcal{O}_X$-modules. It follows from the description of
the functor $\SheafHom(\mathcal{O}_Z, -)$ in
Lemma \ref{lemma-compute-sheaf-with-exact-support}
that there is a map
$$
\Hom_A(B, I^\bullet)
\longrightarrow
\Gamma(Z, \SheafHom(\mathcal{O}_Z, \mathcal{I}^\bullet))
$$
Observe that $\SheafHom(\mathcal{O}_Z, \mathcal{I}^\bullet)$
represents $R\SheafHom(\mathcal{O}_Z, \widetilde{K})$.
Applying the universal property of the $\widetilde{\ }$ functor we
obtain a map
$$
\widetilde{\Hom_A(B, I^\bullet)}
\longrightarrow
R\SheafHom(\mathcal{O}_Z, \widetilde{K})
$$
in $D(\mathcal{O}_Z)$. We may check that this map is an isomorphism in
$D(\mathcal{O}_Z)$ after applying $i_*$. However, once we apply
$i_*$ we obtain the isomorphism of Derived Categories of Schemes,
Lemma \ref{perfect-lemma-quasi-coherence-internal-hom}
via the identification of
Lemma \ref{lemma-sheaf-with-exact-support-ext}.
\end{proof}

\end{document}
